\chapter{The TKS v8 Adapter and the Noetic Topos Boundary}\label{ch09:adapter}

The candidate touches two formal neighbors. Downward, it can hand a single, narrowly qualified class of board component---the certified finite transition---to the TKS v8 finite reference kernel, where the component becomes an exact finite arrow between exact finite objects. Upward, it inherits from the Noetic Topos 2.0 specification an evidence boundary separating \emph{ambient} constructions from \emph{resolved} semantics. Both contacts are boundaries first and bridges second: the adapter is engineered so that crossing it can only ever transport what was already exact, and the topos boundary is inherited so that nothing the candidate produces can drift into semantic authority it has not earned. This chapter specifies the certification conditions under which adaptation is permitted (Section~\ref{ch09:sec:cert}), the treatment of composition and iteration (Section~\ref{ch09:sec:paths}), the adapter certificate field by field (Section~\ref{ch09:sec:certificate}), the rejection of ambient substitution (Section~\ref{ch09:sec:ambient}), the DRC-side face of the same discipline (Section~\ref{ch09:sec:tks001}), the rule that no semantic theorem is created by adaptation (Section~\ref{ch09:sec:notheorem}), the inventory of everything that stays outside the finite arrow and why each exclusion demands its own future types (Section~\ref{ch09:sec:outside}), and finally the ambient/resolved evidence boundary itself (Section~\ref{ch09:sec:topos}).

\section{Certification conditions for adaptation}\label{ch09:sec:cert}

The TKS v8 finite kernel is a small, exact universe: finite objects are named tuples of distinct exact elements, and finite arrows are total functions between them given by explicit index mappings, validated at construction. Nothing about a PCBA board is automatically welcome there. The adapter admits exactly one component kind under exactly one qualification regime.

\begin{definition}[Kernel objects and arrows, as consumed by the adapter]\label{ch09:def:kernel}
For this chapter's purposes, a \emph{finite object} is a pair (name, elements) where the name is a nonempty string and the elements are a finite tuple of pairwise-distinct exact tokens; a \emph{finite arrow} is a quadruple (name, source object, target object, mapping) where the mapping is a tuple of integers, one per source element, each a valid index into the target elements, so that the arrow denotes the total function sending the $i$-th source element to the element at the mapping's $i$-th index. Arrow application is index lookup; nothing else happens at application time. The kernel validates these shape conditions at construction, which is why the adapter can treat a successfully constructed arrow as structurally sound---though it still re-verifies, on its own side, that the constructed arrow's names, carriers, and mapping equal the certified table's (Invariant~\ref{ch09:inv:revalidate}).
\end{definition}

\begin{ruleenv}[Certified transitions only]\label{ch09:rule:certified}
Only a certified \emph{total, finite, deterministic, stateless, non-stochastic} component transition may be adapted to a TKS v8 finite arrow. Concretely, the adapter admits a component if and only if all of the following hold:
\begin{enumerate}
\item The board and component are exact typed values, both revalidated at entry.
\item A fresh DRC run over the board is clean; the adapter refuses any board with errors, and it recomputes the verdict rather than accepting one.
\item The component is present on the board exactly once under its reference, and the board's copy hashes to the same component digest as the value presented---the caller cannot adapt a private variant of a board component.
\item The component's kind is exactly \texttt{finite\_transition} and its semantic profile is exactly \texttt{pcba8.finite-transition.v0}.
\item The component carries no parameters and no state bindings: a finite transition's entire meaning is its table, and undeclared side metadata is refused rather than ignored.
\item The attached transition semantics exist, carry the same profile, and are flagged certified total and deterministic, not stateful, and not stochastic.
\item The transition's certificate digest is fresh: the stored digest equals a recomputation over the canonical table payload, and a payload that cannot be canonically rendered is rejected with the failure quoted.
\end{enumerate}
Each condition failing raises a typed adapter error; there is no partial admission.
\end{ruleenv}

The five adjectives in Rule~\ref{ch09:rule:certified} are the entire theory of the boundary, and each earns its place by what the kernel's arrow type can and cannot represent. \emph{Total}: a finite arrow is defined on every element of its source object; a partial table would require either an error element or an implicit default, both of which are semantic inventions. \emph{Finite}: the source and target carriers must be explicitly enumerated tuples of distinct exact values; nothing with an unbounded or dense carrier can be a kernel object. \emph{Deterministic}: the mapping assigns one target index per source element; a relation or a distribution is not a function. \emph{Stateless}: the arrow's output depends on its input alone; any dependence on history would make the ``function'' a family of functions indexed by hidden state. \emph{Non-stochastic}: the kernel has no randomness anywhere, and a probabilistic transition is a different mathematical object (Section~\ref{ch09:sec:outside}). The adapter checks all five as recorded certificate flags \emph{and} re-derives the checkable ones (totality and index-range validity are re-verified against the actual carriers), so a mislabeled table fails even if its flags lie.

\begin{remark}[Why the DRC is re-run at the adapter]\label{ch09:rem:drcagain}
The DRC already validates finite transitions under \texttt{PCBA8-TKS-001} (Section~\ref{ch09:sec:tks001}), so the adapter's own re-checks are partially redundant by design. The redundancy is the point: the adapter is the last gate before another formal system receives the data, and it treats every upstream assurance as advisory. If the checker and the adapter ever disagree---which should be impossible---the disagreement surfaces as an adapter refusal, never as a kernel object of doubtful pedigree.
\end{remark}

\section{Chains and powers: paths and endomaps}\label{ch09:sec:paths}

Adaptation of single components raises the immediate question of composition: boards chain components, so what does a chain become on the kernel side?

\begin{ruleenv}[Chains become semigroupoid paths]\label{ch09:rule:path}
A typed sequence of certified finite transitions, where each transition's target object matches the next transition's source object exactly, is represented on the kernel side as a \emph{semigroupoid path}: a nonempty, composition-compatible sequence of finite arrows. The path is the honest object: it records the chain as a chain, composable stepwise, without asserting the existence of identities or of any ambient categorical closure beyond what the kernel's finite structure provides.
\end{ruleenv}

\begin{ruleenv}[Powers require a certified endomap]\label{ch09:rule:endomap}
Iteration---applying a transition repeatedly, as declared feedback through a delay would induce tick after tick---is admitted on the kernel side only for a certified \emph{endomap}: a finite arrow whose source and target are the same exact object, element for element and tag for tag. Only an endomap has powers. A transition whose source and target carriers are merely isomorphic, or equal up to renaming, or equal up to retagging of values, is not an endomap and has no powers until the identification is itself made explicit as a certified arrow and composed in.
\end{ruleenv}

Rule~\ref{ch09:rule:endomap} is stricter than casual intuition expects, and deliberately so. ``The same set, basically'' is precisely the reasoning the tagged-value discipline exists to forbid: integer $1$ and string \texttt{"1"} are different elements, an object named \texttt{Stage} and an object named \texttt{Stage'} with equal carriers are different objects, and squaring an arrow between them is a type error, not a formality. The practical consequence for board modeling is clean: a feedback loop's per-tick finite dynamics can be studied as kernel powers exactly when the loop's transition is certified on one fixed carrier---which is also the only situation in which ``the state space of the loop'' is a well-defined finite object rather than a manner of speaking.

\section{The adapter certificate}\label{ch09:sec:certificate}

A successful adaptation returns neither a bare kernel arrow nor a boolean, but an immutable \emph{adapter certificate}: a frozen record binding the arrow to everything that licensed its creation. The certificate is the unit of citation for cross-system claims---any statement of the form ``this board component is represented in the TKS kernel'' is, in the candidate's discipline, a statement that a certificate with a given digest exists.

\begin{definition}[Adapter certificate]\label{ch09:def:certificate}
The certificate carries the following fields, all covered by its own content-derived digest:
\begin{description}
\item[Kernel identity---exact bytes and profile.] The versioned kernel profile string and the SHA-256 digest of the kernel's exact source bytes, computed file by file over the kernel package's module files. This is not a version number; it is the hash of the code that defines the arrow type the certificate holds, pinned to an approved constant. (The constant's 64-hex value is fixed in the implementation at the bound commit; it is not restated in this text, and printing it beside the artifact hashes of Table~\ref{ch03:tab:artifacts} is an open documentation item.) A kernel whose bytes differ---by one character---is not the approved kernel, and adaptation refuses to begin.
\item[Board, component, and transition identities.] The full canonical board payload and the board digest; the full canonical component payload and the component digest; the transition-table certificate digest, recorded twice by design---once as the transition digest and once as the PCBA-side certificate digest---with validation requiring the two to be equal and fresh. The certificate does not point at its board; it contains it.
\item[Tagged carrier.] The exact source and target element tuples, serialized as tagged type-plus-value payloads, together with the exact integer mapping. The carrier travels with the certificate so that the certified function is fully inspectable without dereferencing anything.
\item[Semantic profile and names.] The component's semantic profile string, the component reference (which names the kernel arrow), and the source and target object names.
\item[Fixed boundary.] The Ring-4 research boundary, re-derived and re-checked at every validation, inside the digest as always.
\item[Adapter digest.] The SHA-256 of the certificate's own canonical payload---the identity under which the adaptation is cited.
\end{description}
\end{definition}

\begin{invariant}[Certificate self-verification]\label{ch09:inv:revalidate}
Every public use of a certificate re-runs its full validation, which: revalidates board and component; recomputes both digests and the unique-membership condition of Rule~\ref{ch09:rule:certified}; re-runs the DRC and requires the embedded board to still be clean; reloads the kernel with full provenance checks (Section~\ref{ch09:sec:ambient}) and requires the recorded kernel identity to match the currently loaded one byte for byte; requires the held arrow, its source and target objects, and their element tokens to be exact instances from that kernel, with names, tagged carriers, and mapping all equal to the recorded ones; re-certifies the transition table from the recorded fields and requires the recomputed digest to equal both recorded copies; and finally recomputes the certificate digest over the canonical payload and requires equality. A certificate that fails any clause is not degraded or flagged---it is unusable, and every operation on it raises.
\end{invariant}

Application is likewise disciplined. Calling the certificate on a value performs an exact tagged lookup in the recorded source carrier (a merely equal-looking value of the wrong exact type is outside the carrier and is refused), converts to the kernel's token type, applies the kernel arrow, verifies that the returned token is an exact kernel token belonging to the recorded target carrier, and returns the corresponding recorded target value. The round trip through the kernel is thus checked on both shores: nothing untagged goes in, and nothing unrecognized comes out.

\subsection{Tagged carriers and exact identity}\label{ch09:sub:tagged}

The phrase ``tagged carrier'' compresses the discipline by which the candidate's type-plus-value identity rule crosses the system boundary intact. A finite carrier element is never a bare value; it is a value under a type tag, and equality of elements is equality of (tag, canonical value) pairs. The consequences at the adapter are concrete:

\begin{itemize}
\item The integer $1$, the Boolean true, and the string \texttt{"true"} are three distinct carrier elements. A transition table over $\{$true, false$\}$ is not a table over $\{1, 0\}$, and the adapter will neither confuse them at admission (the distinctness check of Rule~\ref{ch09:rule:certified} operates on tagged keys) nor at application (the lookup of an integer in a Boolean carrier simply fails).
\item Carrier serialization in the certificate is by tagged payloads, so a certificate read back from its canonical form reconstructs exactly the elements that were certified---no locale, encoding, or numeric-widening ambiguity can produce a look-alike carrier.
\item On the kernel side, elements travel as the kernel's own exact token type carrying the same (tag, canonical value) content; the adapter converts on entry, and on exit verifies both the token's exact type and its membership in the recorded target carrier before translating back. A kernel that returned a well-formed token \emph{outside} the certified target carrier---impossible for the approved kernel, but the adapter does not assume approval is competence---would be caught at this membership check.
\end{itemize}

The strictness has one practical corollary for board authors worth stating plainly: two transition components intended to ``share a state space'' share it only if their carriers are element-for-element identical under tags, in the same order, under the same object name. Order matters because the mapping is positional; names matter because kernel objects are named; tags matter because identity is tagged. The candidate provides no carrier-unification convenience, and this is a feature: every identification of state spaces is a modeling decision, and modeling decisions are made on boards, in hashed content, where reviewers can see them.

\section{A worked illustrative adaptation}\label{ch09:sec:worked}

\begin{example}[Illustrative certified transition and its certificate]\label{ch09:ex:worked}
Illustrative only; every identifier and value is constructed for exposition. A board \texttt{demo-fsm} carries a \texttt{finite\_transition} component \texttt{step} whose certified table declares: source object \texttt{Stage} with elements (string \texttt{"idle"}, string \texttt{"primed"}, string \texttt{"fired"}); target object \texttt{Stage} with the same three tagged elements in the same order; and mapping $(1, 2, 2)$---\texttt{idle} goes to \texttt{primed}, \texttt{primed} to \texttt{fired}, \texttt{fired} to \texttt{fired}. The certificate flags assert total and deterministic, deny stateful and stochastic, and the table's digest is fresh. The component has no parameters and no state bindings, and its profile is exactly the finite-transition profile.

Adaptation proceeds: the board revalidates; the DRC re-runs clean (the table also passed \texttt{PCBA8-TKS-001} inside that verdict); \texttt{step} is found exactly once on the board with matching digest; the kernel loads under Protocol~\ref{ch09:pro:kernel}; the three tagged strings become kernel tokens; the kernel constructs the source and target objects and the arrow named \texttt{step} with mapping $(1,2,2)$; and the adapter mints the certificate, whose digest now names this adaptation. Applying the certificate to the string \texttt{"primed"} performs the tagged lookup (index $1$), applies the arrow (index $2$), verifies the returned token, and yields the string \texttt{"fired"}.

Three refusals mark the boundaries. Applying the certificate to the string \texttt{"Primed"} fails: no tagged key matches, and the error states the value is outside the tagged source carrier---no case folding, no fuzzy match. Editing the board (say, adding an unrelated component) and presenting the old certificate fails validation: the embedded board no longer matches the presented lineage, and---because validation re-runs the DRC and recomputes digests---the certificate reports its linkage stale rather than pretending the adaptation still stands. And declaring a variant \texttt{step2} whose flags claim determinism but whose mapping has only two entries fails both the DRC (\texttt{PCBA8-TKS-001}: mapping not total on the source carrier) and, were the DRC somehow bypassed, the adapter's own totality re-derivation.
\end{example}

\begin{example}[Illustrative chain and endomap powers]\label{ch09:ex:powers}
Illustrative only, continuing Example~\ref{ch09:ex:worked}. Because \texttt{step}'s source and target are the same exact object---same name \texttt{Stage}, same tagged elements in the same order---its arrow qualifies as a certified endomap, and kernel powers are defined: $\texttt{step}^2$ has mapping $(2,2,2)$, and $\texttt{step}^n = \texttt{step}^2$ for all $n \geq 2$; the finite dynamics reach a fixed point at \texttt{fired}. Now modify the scenario: a second component \texttt{route} certifies a table from \texttt{Stage} to a two-element object \texttt{Gate} (elements \texttt{"open"}, \texttt{"shut"}). The pair (\texttt{step}, \texttt{route}) is composition-compatible---\texttt{step}'s target equals \texttt{route}'s source exactly---so it forms a semigroupoid path of length two, denoting the composite function; but \texttt{route} has no powers, and no request for them is expressible, because its source and target differ. Finally, suppose an author re-declares \texttt{Stage} in a second table as \texttt{Stage'} with identical elements: the path (\texttt{step}, \texttt{step'}-based arrow) is \emph{not} composition-compatible---the objects differ by name---and the identification, if intended, must itself be certified as an explicit arrow. What looks like pedantry is the discipline of Rule~\ref{ch09:rule:endomap} doing its work: every ``obvious'' identification is either made explicit and checkable or does not exist.
\end{example}

What the two examples jointly license is narrow and exact, per Rule~\ref{ch09:rule:notheorem}: the tables are representable and represented; the powers are properties of the kernel arrow. That the fixed point at \texttt{"fired"} models any process reaching any absorbing state is a semantic claim no artifact in this chapter makes.

\section{Ambient module substitution is rejected}\label{ch09:sec:ambient}

A hash-bound certificate would be worthless if the kernel it names could be quietly swapped at import time---the standard failure being an ambient module on the import path shadowing the intended one, or a monkey-patched class substituted for the arrow type. The adapter treats kernel loading as a provenance problem and solves it with four independent checks.

\begin{protocol}[Exact kernel loading]\label{ch09:pro:kernel}
Before any adaptation or validation, the adapter:
\begin{enumerate}
\item resolves the expected kernel package location relative to its own installed source tree and computes the SHA-256 over the kernel's exact module bytes, file by file with name separators, requiring equality with the approved constant---a kernel edited in place fails here;
\item imports the kernel module and requires the module's reported file path to resolve to exactly the expected location---a same-named module found elsewhere on the path fails here;
\item retrieves the finite-object and finite-arrow types and requires each to be a genuine class whose reported source file resolves to exactly the expected kernel module---a substituted or wrapped class, even one installed into the genuine module's namespace, fails here;
\item records the kernel identity (profile, byte digest, and the package-relative source names) into the certificate, so that later validations repeat the entire procedure and additionally require the recorded identity to equal the freshly derived one.
\end{enumerate}
Any failure raises a typed adapter error naming the untrusted step; no fallback kernel, vendored copy, or degraded mode exists.
\end{protocol}

\begin{threat}[Substitution scenarios the protocol forecloses]\label{ch09:thr:substitution}
Illustrative only, as engineering rationale. (i) A test harness inserts a permissive stub kernel earlier on the import path to ``simplify fixtures''; step 2 rejects the foreign path. (ii) A well-meaning patch adds a convenience subclass of the arrow type and registers it in the kernel's namespace; step 3 rejects the class whose source is not the kernel module, and the exact-type checks of Invariant~\ref{ch09:inv:revalidate} reject instances of any subclass regardless. (iii) A local edit relaxes the kernel's mapping validation; step 1 rejects the changed bytes before the import is even attempted. (iv) A stale certificate from an older kernel is presented after a kernel upgrade; validation's identity-equality clause rejects it, which is the intended behavior---certificates do not survive their kernel, they are re-earned under the new bytes.
\end{threat}

The protocol's philosophy generalizes a theme that runs through the whole candidate: an artifact's dependencies are part of the artifact. The certificate names its kernel by content, and the adapter refuses to reason about arrows whose defining code it cannot re-derive byte for byte.

\section{Certificate lifecycle and staleness}\label{ch09:sec:lifecycle}

Because a certificate re-verifies its entire licensing basis at every use (Invariant~\ref{ch09:inv:revalidate}), it has a well-defined lifecycle governed by one principle: a certificate is valid exactly as long as every identity it binds still holds, and not one moment longer. What does and does not invalidate a certificate follows from the validation clauses rather than from convention, and is enumerated here.

\begin{itemize}
\item \emph{Nothing the caller does to their copy of the board matters}, because the certificate does not reference the caller's board---it embeds its own. Boards are immutable values; a caller who ``edits the board'' has constructed a new board, and the certificate's embedded original is untouched and still valid.
\item \emph{A change to the implementation's source invalidates every certificate}, because validation recomputes the current package identities and the DRC verdict under them; a certificate minted under prior sources fails the freshness comparisons. This is deliberate and symmetric with the simulation result's behavior: artifacts do not outlive the code whose behavior they certify.
\item \emph{A change to the kernel's bytes invalidates every certificate}, through the identity-equality clause of Protocol~\ref{ch09:pro:kernel}. Certificates are re-earned under a new kernel by re-running adaptation, which re-derives everything; there is no migration path that carries an old certificate forward, because a migration would be precisely an unreviewed assertion that the old and new kernels agree.
\item \emph{Time, by itself, changes nothing.} No certificate field carries a timestamp or expiry; the lifecycle is entirely content-driven. Two parties holding the same sources and the same certificate bytes will agree on its validity regardless of when they check.
\end{itemize}

\begin{remark}[Staleness is a feature with a cost, and the cost is accepted]\label{ch09:rem:staleness}
The strict lifecycle means that routine maintenance---reformatting a source file, say---invalidates every outstanding certificate even though ``nothing meaningful changed.'' The candidate accepts this cost knowingly, for the same reason it hashes whole source trees: \emph{meaningful} is exactly the judgment that cannot be delegated to the artifact layer. A digest cannot distinguish a reformatting from a semantics change; a human reviewer can, and the re-adaptation run that refreshes the certificates is cheap, deterministic, and itself an audit event. The alternative---certificates that survive ``insignificant'' changes under some significance policy---would relocate trust from checkable bytes to an uncheckable policy, which is the trade the whole architecture exists to refuse.
\end{remark}

\section{\texttt{PCBA8-TKS-001}: the DRC face of the boundary}\label{ch09:sec:tks001}

The certification conditions of Rule~\ref{ch09:rule:certified} appear twice in the candidate, on independent code paths. Their design-time face is the DRC diagnostic \texttt{PCBA8-TKS-001} (specified in full in Chapter~7): for every \texttt{finite\_transition} component on a board, the checker validates the presence of the table, profile agreement, non-blank object names, pairwise-distinct exact elements in both carriers, totality of the mapping, index-range validity, the certified-flag conjunction (total and deterministic, not stateful, not stochastic), and the freshness and well-formedness of the certificate digest---failing the board on any violation.

The division of labor is deliberate. The DRC catches boundary violations \emph{early}, attached to the offending component, in the same deterministic report as every other board defect, so that an author sees the complete state of the board in one artifact. The adapter catches them \emph{last}, immediately before another formal system is asked to accept the data, trusting nothing upstream. Between the two, the certification conditions are enforced at both ends of the pipeline with no path around either: a board with an invalid table cannot be clean, and a clean board's table is still re-verified at adaptation. Should the two enforcements ever diverge through an implementation defect, the composition fails closed---the adapter's refusal governs.

\section{No automatic semantic theorem}\label{ch09:sec:notheorem}

What does a successful adaptation \emph{prove}? The candidate's answer is deliberately minimal.

\begin{ruleenv}[No-automatic-semantic-theorem rule]\label{ch09:rule:notheorem}
An adapter certificate proves at most this: the named component's transition table is representable, and has been exactly represented, as a type-correct finite arrow in the exact TKS v8 reference kernel identified by the recorded bytes. It carries no automatic CAS semantic theorem whatsoever. In particular, adaptation does not establish that the transition models any cognitive process; that the source or target carriers correspond to any TKS object, world, foundation, or element beyond the names the board author chose; that composition of adapted arrows mirrors any composition of cognitive operations; or that any property provable of the arrow in the kernel (totality, functionality, iterability of an endomap) transfers to any system outside the software. Semantic claims connecting kernel arrows to CAS or TKS meaning require their own separately stated, separately reviewed arguments, none of which is supplied by the adapter.
\end{ruleenv}

The rule guards against a subtle laundering pattern: formal systems confer prestige, and an object that has ``made it into the kernel'' can appear validated by association. The candidate's semantic-identity discipline (namespaced references, no silent identity across namespaces) already blocks the crudest version---two things sharing a label do not thereby share a referent---and Rule~\ref{ch09:rule:notheorem} blocks the refined version: two things sharing a \emph{formal representation} do not thereby share a semantics. The adapter is a typed transport, not an oracle; what it transports faithfully is structure, and structure is all it certifies.

Three misreadings recur in adjacent literatures; each is refuted here, stated as the inference it tempts and the reason the inference fails.

\emph{``The kernel accepted it, so the model is coherent.''} The kernel accepted a total function between finite tagged sets. Coherence of the \emph{modeling}---whether the carrier elements carve any phenomenon at joints, whether the mapping reflects any dynamics---was never examined by anything in the pipeline. A table mapping arbitrary strings arbitrarily adapts exactly as smoothly as a carefully theorized one; the certificate cannot tell them apart, and does not claim to.

\emph{``The arrow composes/iterates cleanly, so the modeled process chains/repeats.''} Composition and powers are closure properties of the representation. The modeled process, if any, has whatever closure properties it has; the board author's claim that stepwise application models repeated exposure, or that composition models sequential processing, is a hypothesis about the world stated in the modeling layer, and the kernel's algebra neither strengthens nor tests it.

\emph{``The carrier elements are named after TKS objects, so the adapter connects CAS to TKS semantics.''} Names are safe references, not imports. The semantic-identity tuple that would bind a carrier element to an authoritative TKS definition (namespace, local identity, version, source identity, definition digest) is not part of the transition table, and no adapter field carries it. A board author who names an element \texttt{"B5"} has named a string. The candidate's separation between structural reference and semantic identity is precisely what makes this safe: nothing downstream will ever mistake the string for the object, because nothing downstream resolves strings to objects at all.

The positive statement hiding under the three refutations is worth making once: what the adapter \emph{does} establish---exact structural transport with content-derived accountability---is the necessary substrate for any future semantic argument. A semantic theorem connecting a board's finite dynamics to a TKS-object-level claim would need, at minimum, an exact and stable representation on both sides to even state its correspondence. The adapter supplies the PCBA side of that substrate and nothing more; the theorem, if one is ever proposed, will be a new artifact with its own review.

\section{What stays outside the finite arrow}\label{ch09:sec:outside}

The finite arrow's admission conditions exclude most of the candidate's own vocabulary, and the exclusions are permanent features of the boundary, not backlog. Each excluded class fails a different admission condition, and each therefore requires a different future type with different comparison maps; coercing any of them into finite arrows would produce objects that type-check and lie.

\begin{boundary}\label{ch09:bd:outside}
The following remain outside the finite arrow, each for the stated structural reason:
\begin{description}
\item[Continuous simplex profiles.] A strictly positive normalized profile over the forty ND-cells inhabits an uncountable open simplex. No enumeration of its carrier exists, so it can never be a finite object's element list. Its natural home is a metric/geometric type (the mixed-simplex profiles of Chapter~6 with their disclosed numerical limits); relating it to finite data requires explicit discretization maps whose loss is declared, not a coercion that pretends the simplex was finite.
\item[Response curves.] A component transfer rule such as the junction's clipped product is a function on a dense decimal domain. Sampling it at finitely many points yields a finite table \emph{about} the curve, not the curve; the table's fidelity is a claim needing its own error bounds. A future curve type must carry domain, arithmetic context, and approximation contract---three things a finite arrow has no slot for.
\item[Instrumented observations.] A measurement record is not a value; it is a value plus provenance: tick, producer, calibration reference, uncertainty, measurement kind, scope digest. Stripping the provenance to admit the value into a carrier would be exactly the laundering the measurement schema exists to prevent. Observation types belong to the evidence layer, where their calibration metadata is load-bearing (Section~\ref{ch09:sec:topos}).
\item[Stateful devices.] The delay and the latching fuse compute functions of history, not of input. The honest formalization of a stateful device is a transition \emph{system}---state carrier, update map, output map---and while each component map might individually be finite, the device is the system, not either map. A future stateful profile must type the state explicitly; the finite arrow's statelessness condition exists precisely to refuse the alternative, an arrow with hidden memory.
\item[Stochastic kernels.] A probabilistic transition assigns each source element a distribution over targets. That is a Markov kernel, a different object with different composition (distributions compose by integration, not by function application) and different equality (almost-sure, not pointwise). Determinism is an admission condition because a finite arrow asserting one output per input would simply misreport a stochastic table.
\item[Memristive history variables.] Device semantics whose parameters drift with accumulated exposure---the memristive elements the source book gestures at---are path-dependent: their behavior is a function of an unbounded history summary. They compound the stateful and continuous exclusions and would additionally require declared time-integration semantics, which the candidate's synchronous tick discipline does not supply. They are named in the open-obligations register precisely so that their absence is a recorded boundary, not an oversight.
\end{description}
In every case the future path is the same: a distinct type, a declared comparison map to and from existing types, and validation of that map---never coercion into the finite arrow.
\end{boundary}

\subsection{Comparison maps, not coercions}\label{ch09:sub:comparison}

The closing sentence of Boundary~\ref{ch09:bd:outside}---distinct types with declared comparison maps, never coercion---states the general method for every future extension, and its content is best seen by contrasting the two disciplines on a single case. Take the response curve. A \emph{coercion} discipline would say: sample the junction's clipped product at, say, eleven points, build a finite table, and adapt it; the table type-checks, and downstream consumers cannot tell it from a genuinely finite transition. A \emph{comparison-map} discipline says instead: define a curve type carrying the rule, domain, and arithmetic context; define an explicit sampling map from curves to finite tables carrying its sample grid as data; and let every consumer of the resulting table see, in the types, that it holds a \emph{sampled image of a curve} with a named grid---an object about which the honest questions (What happens between grid points? Was the grid adequate for the use?) remain askable because the provenance was not erased.

The difference is not pedantic. Under coercion, the information that a finite object approximates a continuous one is destroyed at the boundary, and every downstream inference silently inherits an undeclared adequacy assumption. Under comparison maps, the same inference must cite the map, and the map's adequacy is a reviewable claim with an owner. The candidate applies the identical pattern to each excluded class: discretization maps for simplex profiles (with declared loss), sampling maps for curves (with declared grids), summarization maps for histories (with declared state abstractions), support restrictions or derandomizations for stochastic kernels (with declared policies)---each a future, versioned, tested artifact. None exists at the bound sources; the pattern is recorded so that when they are built, they are built as maps and not as leaks.

\begin{remark}[The boundary as a design asset]\label{ch09:rem:asset}
It may seem that so restrictive an adapter is barely useful. The restriction is the utility. Because admission is so narrow, the certificate's guarantee is total over its domain: every certified arrow is exactly faithful, with no approximation regime, no sampling caveat, no hidden state. The excluded classes are not banished from the candidate---they live in their own layers with their own contracts. The boundary merely ensures that the one bridge between systems carries only cargo that survives the crossing without loss.
\end{remark}

\section{The Noetic Topos evidence boundary}\label{ch09:sec:topos}

The candidate's second formal contact is read-only. From the Noetic Topos 2.0 specification it inherits a single discipline governing how raw records relate to semantic conclusions: the distinction between the \emph{ambient} construction, where records may be associated as external evidence, and \emph{resolved} semantics, which records may enter only through defined and tested restriction and gluing operations.

\begin{ruleenv}[Ambient association only]\label{ch09:rule:ambient}
Raw or partially resolved records---simulation measurements, test-point observations, locally calibrated traces, adapter certificates---may be treated as external evidence \emph{associated with} the ambient construction. Association is bookkeeping: it records that evidence exists and where it attaches. It confers no semantic status. In particular, a \emph{marked cover proves no measurement true}: marking a family of observation sites as jointly covering a region of interest is a declaration of intended coverage, and neither the marking, nor the cover's completeness, nor agreement among the covered observations converts any observation into a true, resolved semantic fact.
\end{ruleenv}

\begin{ruleenv}[Resolution requires restriction and gluing under one calibration profile]\label{ch09:rule:gluing}
No record enters resolved semantics until restriction and gluing are \emph{defined and tested} for the record class under \emph{one} calibration profile: a declared operation restricting a record to a sub-context, a declared operation gluing compatible records on overlapping contexts into one record on the union, declared compatibility conditions on overlaps, and tests exercising all three---all relative to a single, versioned calibration profile, so that ``compatible'' never silently means ``measured under different regimes that happen to agree numerically.'' At the bound sources, no such operations have been defined or tested for any PCBA record class. Consequently the candidate holds \emph{no} resolved semantics: every measurement, trace, and certificate in this monograph is ambient evidence, in full.
\end{ruleenv}

\begin{example}[Illustrative failure of naive gluing]\label{ch09:ex:gluing}
Illustrative only. Suppose two simulation results over the same board region exist: run $A$ observed the effect net through test point \texttt{tp1} under calibration profile \texttt{cal-v1}, and run $B$ observed the same net under a revised profile \texttt{cal-v2} whose declared uncertainty and protocol identity differ. Both artifacts are internally valid; on the overlapping tick range their observed value texts happen to be numerically equal. A naive gluing rule---``records over the same net that agree on overlaps may be merged into one observation series''---would accept the merge, and the merged series would then appear to be a single, longer observation of the net under \emph{a} calibration, when in fact no single calibration covers it and the numerical agreement is uninterpreted coincidence as far as the metadata is concerned. Rule~\ref{ch09:rule:gluing} refuses the merge twice over: no gluing operation has been defined or tested for observation records at all, and even a future operation would be scoped to one calibration profile, which this pair fails by construction. The honest artifact remains two results, each citable in its own scope---and the example shows why the single-profile condition is not bureaucratic caution but the difference between an evidence series and a concatenation of look-alikes.
\end{example}

The sheaf-theoretic intuition behind the rule can be made explicit without claiming any theorem. Gluing is where local observations become a global assertion---it is exactly the step at which evidence is promoted to fact. A formalism that lets records glue by default therefore manufactures global facts from local bookkeeping, and does so most confidently where instrumentation is least consistent. The topos boundary places the promotion step behind explicit definitions and tests, and the single-calibration-profile condition pins the one variable that most easily invalidates a glue: two locally plausible observations under different calibration regimes may satisfy any numerical compatibility check and still describe incomparable quantities. Until the operations exist, are versioned, and are tested, the honest global statement available from any family of PCBA records is the conjunction of its members' individual, licensed claims---tested software behavior, exactly bound---and nothing stronger.

\begin{boundary}\label{ch09:bd:closing}
Summary of the chapter's licensed claims. The adapter transports certified total, finite, deterministic, stateless, non-stochastic transitions---and nothing else---into the exact TKS v8 kernel, returning immutable certificates that bind kernel bytes, board, component, transition, tagged carrier, semantic profile, and the fixed Ring-4 boundary, and that re-verify all of it at every use. Chains are semigroupoid paths; powers exist only for certified endomaps; ambient kernel substitution is structurally rejected; \texttt{PCBA8-TKS-001} enforces the same conditions at design time. A certificate proves type-correct representation and nothing more: no CAS semantic theorem, no cognitive correspondence, no transfer of kernel-provable properties to any system outside the software. Continuous, curved, instrumented, stateful, stochastic, and history-dependent semantics remain outside the finite arrow pending their own types and comparison maps. All records are ambient evidence under the Noetic Topos boundary; a marked cover proves no measurement true; and no resolved semantics exist, or are claimed, until restriction and gluing under one calibration profile are defined and tested. All artifacts remain Ring~4, \textsc{not licensed}, and every obligation named here is open.
\end{boundary}
